\documentclass[11pt,a4paper]{article}
\usepackage[margin=2.5cm]{geometry}
\usepackage{booktabs}
\usepackage{tabularx}
\usepackage{multirow}
\usepackage{xcolor}
\usepackage{hyperref}
\usepackage{enumitem}
\usepackage{array}
\usepackage{colortbl}
\usepackage{parskip}

\hypersetup{
    colorlinks=true,
    linkcolor=blue!60!black,
    urlcolor=blue!60!black,
}

\definecolor{planA}{HTML}{D6EAF8}
\definecolor{planB}{HTML}{D5F5E3}
\definecolor{planC}{HTML}{FDEBD0}
\definecolor{covered}{HTML}{27AE60}
\definecolor{notcovered}{HTML}{E74C3C}
\definecolor{headerrow}{HTML}{2C3E50}

\newcommand{\rev}[1]{\texttt{\textbf{#1}}}
\newcommand{\expid}[1]{\textbf{#1}}
\newcommand{\yes}{\textcolor{covered}{$\bullet$}}
\newcommand{\no}{\textcolor{notcovered}{--}}

\title{\textbf{SAGE ICML 2026 Rebuttal}\\[0.4em]
\large Experiment Plan: Minimal, Average, and Maximal Tiers}
\author{}
\date{March 2026}

\begin{document}
\maketitle
\tableofcontents
\newpage

% =========================================================
\section{Paper and Review Summary}
% =========================================================

\textbf{SAGE} (\textit{Self-Assessed Gradient-based Enhancement}) proposes test-time alignment
without retraining or external reward models. Key components: (1) the base LLM acts as a
soft self-judge via calibrated contrastive log-probability margins; (2) candidates are partitioned
into good/bad groups, yielding a stable \textit{textual gradient}; (3) generation is iteratively
refined conditioned on this gradient.

\textbf{Current results (main table):} Qwen3-8B and Llama3-8B-Instruct on MATH-500 (92.0\%),
AlpacaEval (win rate 74.9), and XSTest (93.5\%).

\subsection{Reviewer Scores}

\begin{center}
\begin{tabular}{lllll}
\toprule
\textbf{Reviewer} & \textbf{ID} & \textbf{Recommendation} & \textbf{Confidence} & \textbf{Soundness} \\
\midrule
R1 & \rev{YbZE} & 3 --- Weak reject & 2 & 2 \\
R2 & \rev{aHH9} & 3 --- Weak reject & 3 & 3 \\
R3 & \rev{RD4w} & 4 --- Weak accept & 3 & 3 \\
R4 & \rev{pDvu} & 3 --- Weak reject & 3 & 3 \\
\bottomrule
\end{tabular}
\end{center}

% =========================================================
\section{Reviewer Concern Map}
% =========================================================

\begin{table}[h!]
\centering
\renewcommand{\arraystretch}{1.3}
\begin{tabularx}{\textwidth}{llX}
\toprule
\textbf{Code} & \textbf{Reviewer} & \textbf{Concern} \\
\midrule
YbZE-Q1 & \rev{YbZE} & Small models ($<$3B): does the textual gradient become noise? \\
YbZE-Q2 & \rev{YbZE} & Wall-clock time comparison with Best-of-N under same end-to-end latency \\
\midrule
aHH9-Q1 & \rev{aHH9} & Wall-clock time for full SAGE run vs.\ Best-of-N+RM \\
aHH9-Q2 & \rev{aHH9} & Tested on larger models (70B+)? Blind spots at small scale? \\
aHH9-W4 & \rev{aHH9} & AlpacaEval gains may be verbosity-driven, not quality \\
aHH9-Q3 & \rev{aHH9} & Can SAGE get stuck in a ``self-congratulatory'' loop reinforcing wrong paths? \\
\midrule
RD4w-Q1 & \rev{RD4w} & Why was GPT-4.1 used for SPO prompt tuning? Unfair baseline? \\
RD4w-Q2 & \rev{RD4w} & How are aspects determined? Sensitivity to different aspect sets? \\
RD4w-Q3 & \rev{RD4w} & How should $m_{\min}$ (group size) be selected in practice? \\
RD4w-Q4 & \rev{RD4w} & Cost and time differences between SAGE and baselines? \\
\midrule
pDvu-W1 & \rev{pDvu} & Qwen3-8B baseline discrepancy: paper reports 84.4, tech report gives 87.4 (non-thinking) / 97.4 (thinking). Gains may not hold in reasoning mode. \\
pDvu-W2 & \rev{pDvu} & Benchmarks too easy for Qwen3-8B (already at 84\%+); harder benchmarks needed. \\
pDvu-W3 & \rev{pDvu} & RM baseline (\texttt{sfairXC/FsfairX-LLaMA3-RM-v0.1}) is outdated vs.\ Qwen3-8B. \\
\bottomrule
\end{tabularx}
\caption{Mapped reviewer concerns (13 total across 4 reviewers).}
\end{table}

% =========================================================
\section{Recommended Execution Order}
% =========================================================

\begin{center}
\begin{tabular}{lll}
\toprule
\textbf{Day(s)} & \textbf{Experiment} & \textbf{Note} \\
\midrule
1     & A2 (Latency Table)            & Fastest; answers 4 questions at once \\
1--2  & A1 (Baseline Clarification)   & Runs in parallel with A2 \\
2--3  & A3 (IFEval)                   & Core new benchmark \\
\multicolumn{3}{l}{\quad $\Rightarrow$ \textit{Plan A complete; rebuttal can be submitted here}} \\
\midrule
4--5  & B2 ($m_{\min}$ Ablation)      & 5 configs; parallelizable \\
5--6  & B3 (Aspect Sensitivity)       & 3 configs $\times$ 2 benchmarks \\
6--7  & B4 (Small Model)              & Fast at 1.7B \\
\multicolumn{3}{l}{\quad $\Rightarrow$ \textit{Plan B complete}} \\
\midrule
8--11 & C1 (Qwen3-32B)                & Long; use 2 GPUs \\
11--12& C2 (Updated RM)               & Moderate cost \\
12--14& C3 (MMLU-Pro STEM)            & Broad coverage \\
\multicolumn{3}{l}{\quad $\Rightarrow$ \textit{Plan C complete}} \\
\bottomrule
\end{tabular}
\end{center}

% =========================================================
\section{Plan A --- Minimal ($\sim$10 H100-hours, $\sim$3 days)}
% =========================================================

\textbf{Goal:} Address the three highest-risk concerns to flip weak rejects. Submit a rebuttal
at day 3 if under time pressure.

\subsection*{A1 --- Baseline Clarification and Qwen3-8B Reasoning Mode Reference}

\begin{description}[leftmargin=2cm, style=nextline]
  \item[Covers] \rev{pDvu}-W1
  \item[Rationale] The 84.4 baseline in the paper corresponds to Qwen3-8B in \emph{non-thinking}
    mode with the paper's exact prompt format. The tech report's 87.4 uses a different system
    prompt configuration; the 97.4 figure uses \emph{thinking} mode, which is a separate inference
    regime with a much higher compute budget. These are not directly comparable to SAGE's
    non-thinking setup. No new SAGE run is required for this: add a single reference row in
    Table~1 showing Qwen3-8B in thinking mode ($\approx$97.4) and frame SAGE as a
    complementary method that operates in the non-thinking regime and closes the gap.
  \item[Model] Qwen3-8B (thinking mode, reference only)
  \item[Benchmark] MATH-500
  \item[GPU estimate] $\sim$1 H100-hour
  \item[Wall-clock] 0.5 days
  \item[Expected outcome] Confirms 84.4 is the correct non-thinking baseline. Thinking mode
    reaches $\approx$97 without SAGE. SAGE at 92.0 is a strong result within the non-thinking
    budget. Framing defuses \rev{pDvu}-W1 without requiring SAGE to beat the reasoning mode.
\end{description}

\subsection*{A2 --- Wall-Clock Latency Table}

\begin{description}[leftmargin=2cm, style=nextline]
  \item[Covers] \rev{YbZE}-Q2, \rev{aHH9}-Q1, \rev{aHH9}-W1, \rev{RD4w}-Q4
    \emph{(one table answers four reviewer questions)}
  \item[Procedure] Run SAGE, Best-of-N (same $N$), and Best-of-N+RM on 100 MATH-500 problems.
    Measure: (a) GPU seconds/problem and (b) wall-clock seconds/problem at \texttt{batch\_size=1}
    simulating a single user. Also compute accuracy/time as a composite efficiency metric.
  \item[Model] Qwen3-8B
  \item[Benchmark] MATH-500 (100-problem subset sufficient for timing statistics)
  \item[GPU estimate] $\sim$2 H100-hours
  \item[Wall-clock] 0.5 days
  \item[Expected outcome] SAGE is slower per problem in absolute wall-clock but achieves higher
    accuracy per unit of inference compute once the judge pass is amortized over a batch.
    A Pareto-efficiency plot (accuracy vs.\ wall-clock) shows SAGE dominates at moderate latency
    budgets.
\end{description}

\subsection*{A3 --- IFEval (Verifiable, Judge-Free Instruction Following)}

\begin{description}[leftmargin=2cm, style=nextline]
  \item[Covers] \rev{aHH9}-W4, \rev{pDvu}-W2
  \item[Rationale] IFEval contains 541 prompts across 25 instruction categories (length
    constraints, format rules, keyword presence/absence). Correctness is evaluated by a
    deterministic Python script --- no GPT judge, no verbosity bias. Qwen3-8B scores
    $\approx$78--82\% prompt-level accuracy out-of-the-box, leaving meaningful headroom.
    Crucially: if SAGE only inflates verbosity, it will \emph{fail} IFEval's
    length-limit instructions (e.g., ``respond in under 100 words''), directly falsifying
    \rev{aHH9}-W4. This is the single most informative benchmark to add.
  \item[Model] Qwen3-8B
  \item[Benchmark] IFEval (\texttt{google/IFEval}), evaluated with the official script
  \item[GPU estimate] SAGE: $\sim$4 H100-hours; baseline + Best-of-N: $\sim$1 H100-hour
  \item[Wall-clock] 1 day
  \item[Expected outcome] SAGE improves prompt-level accuracy by 3--7 points over baseline.
    Improvement on length-constrained instructions directly refutes the verbosity hypothesis.
    Also answers \rev{pDvu}-W2 by demonstrating gains on a benchmark where the base model
    has substantial room to improve.
\end{description}

\paragraph{Plan A totals.} $\approx$7--10 H100-hours, 2--3 days. Covers \textbf{7 of 13} concerns.

% =========================================================
\section{Plan B --- Average ($\sim$40 H100-hours, $\sim$7--10 days)}
% =========================================================

\textbf{Goal:} Add breadth and fill the missing ablations that \rev{RD4w} (weak accept) and
\rev{YbZE} explicitly requested. Covers nearly all concerns.

Plan B includes all experiments from Plan~A plus the following three experiments.

\subsection*{B2 --- Group Size ($m_{\min}$) Ablation}

\begin{description}[leftmargin=2cm, style=nextline]
  \item[Covers] \rev{RD4w}-W3, \rev{RD4w}-Q3
  \item[Procedure] Sweep $m_{\min} \in \{1, 2, 4, 6, 8\}$ on MATH-500 (full 500 problems), with
    all other hyperparameters fixed. $m_{\min}=1$ is the degenerate best-vs.-worst pairwise
    case that \rev{RD4w} specifically asked about. Report a figure with $m_{\min}$ on the x-axis
    and accuracy on the y-axis, plus a practical selection heuristic derived from the curve.
  \item[Model] Qwen3-8B
  \item[Benchmark] MATH-500
  \item[GPU estimate] 5 configs $\times$ $\approx$2 H100-hours $=$ \textbf{10 H100-hours}
    (parallelizable across 5 GPUs to 0.5 days)
  \item[Wall-clock] 2 days (serial), 0.5 days (parallel)
  \item[Expected outcome] Performance degrades at $m_{\min}=1$ (noisy single-pair gradient),
    peaks around $m_{\min}=4$--6, and plateaus. Confirms that grouped selection is strictly
    better than best-vs.-worst and that $m_{\min}=4$ is a safe default. Converts
    \rev{RD4w}-W3 from ``no proof'' to ``see Figure~X.''
\end{description}

\subsection*{B3 --- Aspect Sensitivity Ablation}

\begin{description}[leftmargin=2cm, style=nextline]
  \item[Covers] \rev{RD4w}-W2, \rev{RD4w}-Q2
  \item[Procedure] Run three aspect configurations: (a)~paper default aspects,
    (b)~single generic aspect (``overall quality and correctness''), (c)~task-specific
    aspects for math (logical coherence, calculation accuracy, step completeness).
    Evaluate on MATH-500 and IFEval. Report a $3 \times 2$ table.
  \item[Model] Qwen3-8B
  \item[Benchmarks] MATH-500, IFEval
  \item[GPU estimate] 3 configs $\times$ 2 benchmarks $\times$ $\approx$2 H100-hours
    $=$ \textbf{12 H100-hours}
  \item[Wall-clock] 2 days
  \item[Expected outcome] Default aspects outperform the single generic aspect by a small but
    consistent margin. Task-specific aspects match or marginally exceed the defaults.
    Demonstrates that aspect choice matters but the method is not brittle to exact
    formulation, directly answering \rev{RD4w}-Q2.
\end{description}

\subsection*{B4 --- Small Model Test (Qwen3-1.7B)}

\begin{description}[leftmargin=2cm, style=nextline]
  \item[Covers] \rev{YbZE}-Q1, \rev{YbZE}-W1
  \item[Rationale] Qwen3-1.7B is below the 3B threshold \rev{YbZE} cited. It has sufficient
    instruction-following for the failure mode (if any) to be informative rather than
    trivially broken.
  \item[Procedure] Run SAGE with Qwen3-1.7B on MATH-500 and IFEval. Also track the variance
    of log-probability margins across iterations as a noise proxy. Compare against Qwen3-1.7B
    baseline and Best-of-N.
  \item[Model] Qwen3-1.7B (Qwen2.5-3B as fallback if 1.7B proves too unstable)
  \item[Benchmarks] MATH-500, IFEval
  \item[GPU estimate] $\approx$5 H100-hours (fast at 1.7B)
  \item[Wall-clock] 1 day
  \item[Expected outcome] Modest gains on IFEval (lower-complexity instructions),
    smaller or no gains on MATH-500. The log-probability margin variance plot shows
    when noise becomes dominant. Provides an honest empirical boundary condition:
    SAGE is less effective below 3B on hard reasoning but not destructively noisy ---
    directly answering \rev{YbZE}-Q1.
\end{description}

\paragraph{Plan B totals.} $\approx$38--40 H100-hours, 7--10 days. Covers \textbf{10 of 13} concerns.

% =========================================================
\section{Plan C --- Maximal ($\sim$80 H100-hours, $\sim$20--25 days)}
% =========================================================

\textbf{Goal:} Exhaustive coverage. Warranted only if reviewers raise scores after Plan~B,
or if you have $>$2 weeks and want a stronger camera-ready paper.

Plan C includes all experiments from Plans A and B plus the following three experiments.

\subsection*{C1 --- Qwen3-32B}

\begin{description}[leftmargin=2cm, style=nextline]
  \item[Covers] \rev{aHH9}-Q2 (partial), \rev{RD4w}-W1 (partial)
  \item[Rationale] Qwen3-32B is meaningfully stronger than 8B and fits on 2$\times$H100~80GB
    with tensor parallelism or standard 8-bit quantization. Testing 70B+ (as \rev{aHH9} asked)
    would require 4$\times$H100 for marginal reviewer satisfaction; 32B is the best ROI point.
  \item[Procedure] Run SAGE and baseline Qwen3-32B on MATH-500 and IFEval. Compare against
    Qwen3-32B Best-of-N. If SAGE still improves over a strong 32B baseline, the argument
    that ``small model blind spots undermine self-assessment'' is substantially weakened.
  \item[Model] Qwen3-32B (2$\times$H100, tensor parallel)
  \item[Benchmarks] MATH-500, IFEval
  \item[GPU estimate] $\approx$20 H100-hours
  \item[Wall-clock] 3--4 days
  \item[Expected outcome] SAGE adds +2--4 points on MATH-500 and +3--5 points on IFEval for
    Qwen3-32B. The per-step gain is smaller than at 8B (stronger base $=$ less headroom) but
    the absolute accuracy is higher, demonstrating SAGE scales with model capability rather
    than depending on model weakness.
\end{description}

\subsection*{C2 --- Updated Reward Model Baseline}

\begin{description}[leftmargin=2cm, style=nextline]
  \item[Covers] \rev{pDvu}-W3
  \item[Procedure] Replace \texttt{sfairXC/FsfairX-LLaMA3-RM-v0.1} with a Qwen3-era reward
    model (candidates: \texttt{Skywork-Reward-Qwen-2.5-7B-v0.2} or
    \texttt{ArmoRM-Llama3-8B-v0.1} as fallback). Re-run the Best-of-N+RM comparison on
    MATH-500 and IFEval.
  \item[Model] Qwen3-8B + updated RM
  \item[Benchmarks] MATH-500, IFEval
  \item[GPU estimate] $\approx$8 H100-hours
  \item[Wall-clock] 1.5 days
  \item[Expected outcome] If SAGE matches or exceeds the updated RM: the objection is fully
    resolved. If SAGE slightly trails: frame honestly as ``SAGE is a single-model solution
    that approaches the performance of a two-model pipeline without the deployment overhead,''
    which is a legitimate practical advantage.
\end{description}

\subsection*{C3 --- MMLU-Pro STEM Subset}

\begin{description}[leftmargin=2cm, style=nextline]
  \item[Covers] \rev{pDvu}-W2 (secondary), \rev{RD4w}-W1 (secondary)
  \item[Rationale] MMLU-Pro STEM contains $\approx$2000 questions. Qwen3-8B scores
    $\approx$55--62\% out-of-the-box --- substantial headroom. Multiple-choice
    evaluation is fully verifiable without a judge, and the domain breadth (biology,
    chemistry, physics, engineering) extends beyond SAGE's current math-only evaluation.
  \item[Model] Qwen3-8B
  \item[Benchmark] MMLU-Pro STEM split
  \item[GPU estimate] $\approx$12 H100-hours
  \item[Wall-clock] 2 days
  \item[Expected outcome] SAGE improves by +3--6 points over baseline. This demonstrates
    that gains are not specific to math and are not caused by verbosity inflation,
    complementing the IFEval result.
\end{description}

\paragraph{Plan C totals.} $\approx$78--80 H100-hours, 20--25 days. Covers \textbf{all 13} concerns.

% =========================================================
\section{Concern Coverage by Plan}
% =========================================================

\begin{table}[h!]
\centering
\renewcommand{\arraystretch}{1.4}
\begin{tabular}{llccc}
\toprule
\textbf{Code} & \textbf{Concern (short)} & \textbf{Plan A} & \textbf{Plan B} & \textbf{Plan C} \\
\midrule
YbZE-Q1 & Small model $<$3B noise         & \no  & \yes\ B4 & \yes\ B4 \\
YbZE-Q2 & Wall-clock vs.\ BoN             & \yes\ A2 & \yes\ A2 & \yes\ A2 \\
aHH9-Q1 & Wall-clock vs.\ BoN+RM          & \yes\ A2 & \yes\ A2 & \yes\ A2+C2 \\
aHH9-Q2 & Tested on 70B+?                 & \no  & \no  & \yes\ C1 (32B) \\
aHH9-W4 & AlpacaEval verbosity            & \yes\ A3 & \yes\ A3 & \yes\ A3 \\
aHH9-Q3 & Self-congratulatory loop        & \no  & \no  & \no \\
RD4w-Q1 & GPT-4.1 for SPO (text response) & \yes* & \yes* & \yes* \\
RD4w-Q2 & Aspect determination            & \no  & \yes\ B3 & \yes\ B3 \\
RD4w-Q3 & $m_{\min}$ selection            & \no  & \yes\ B2 & \yes\ B2 \\
RD4w-Q4 & Cost/time vs.\ baselines        & \yes\ A2 & \yes\ A2 & \yes\ A2 \\
pDvu-W1 & Baseline discrepancy            & \yes\ A1 & \yes\ A1 & \yes\ A1 \\
pDvu-W2 & Benchmarks too easy             & \yes\ A3 & \yes\ A3 & \yes\ A3+C3 \\
pDvu-W3 & Outdated RM                     & \no  & \no  & \yes\ C2 \\
\midrule
\multicolumn{2}{l}{\textbf{Total covered}} & \textbf{7/13} & \textbf{10/13} & \textbf{12/13} \\
\bottomrule
\end{tabular}
\caption{Coverage of reviewer concerns per plan tier. \yes\ = covered with experiment ID.
\no\ = not covered. * = addressed via rebuttal text, no experiment needed.}
\end{table}

% =========================================================
\section{Rebuttal Text Notes (No Experiments Needed)}
% =========================================================

\begin{description}[leftmargin=2cm, style=nextline]
  \item[\rev{pDvu}-W1 (Baseline discrepancy)]
    Explicitly state in the rebuttal that 84.4 matches the Qwen3-8B non-thinking baseline
    under the paper's exact prompt format. The tech report's 87.4 uses a different system
    prompt (cite the evaluation configuration section of the Qwen3 technical report). The 97.4
    figure is the thinking (reasoning) mode result with a substantially larger compute budget ---
    a separate inference regime. SAGE and the thinking mode are not in competition; they operate
    under different compute regimes.

  \item[\rev{RD4w}-Q1 (GPT-4.1 for SPO)]
    SPO requires an external optimizer by design; using a weaker model for SPO prompt tuning
    would hurt SPO further, making the comparison even more conservative for SAGE. This choice
    follows the original SPO paper's setup. SAGE requires no external model whatsoever.

  \item[\rev{aHH9}-Q3 (Self-congratulatory loop)]
    The loop concern is partially addressed by Figure~4(b) in the paper, which shows the RM
    score dips before recovering --- indicating the self-judge does \emph{not} monotonically
    reinforce the same path. The textual gradient is regenerated each iteration from newly
    grouped candidates, preventing static reinforcement. An additional analysis plotting
    response diversity (e.g., embedding-space variance across epochs) from Figure~11 of the
    Appendix can be cited here.

  \item[\rev{aHH9}-W3 (Reflection ability varies)]
    Reframe as expected behavior: the method's effectiveness naturally scales with the model's
    reflective capacity, which is why Qwen3-8B shows larger absolute gains than Llama3-8B-Instruct.
    Experiment B4 (Qwen3-1.7B) provides the empirical lower bound on model scale.
\end{description}

\end{document}
